\documentclass[11pt]{article}

\usepackage{amsmath,amssymb,mathtools,bm}
\usepackage[margin=1in]{geometry}
\usepackage{setspace}

\newcommand{\E}{\mathbb{E}}
\newcommand{\R}{\mathbb{R}}
\newcommand{\N}{\mathcal{N}}
\newcommand{\ELBO}{\mathrm{ELBO}}
\newcommand{\tr}{\mathrm{tr}}
\setlength{\parindent}{0pt}

\title{20260401A Number Crunch}
\author{}
\date{}

\begin{document}
\maketitle

\section*{Background and Problem Intuition}
\textbf{Mathematical process:} Start from Bayes' rule, then convert inference into optimisation.
\begin{align}
p(\bm\theta\mid\bm y)
&= \frac{p(\bm y\mid\bm\theta)p(\bm\theta)}{p(\bm y)}
 = \frac{p(\bm y\mid\bm\theta)p(\bm\theta)}{\displaystyle\int p(\bm y,\bm\theta)\,d\bm\theta} \\ 
p(\bm y)
&= \int p(\bm y,\bm\theta)\,d\bm\theta\\
&= \int p(\bm y \mid \bm\theta) p(\bm\theta)\,d\bm\theta, \\[1em] 
\intertext{{\setstretch{1.5}where $\displaystyle\int p(\bm y\mid\bm\theta)\,p(\bm\theta)\,d\bm\theta$ is potentially
intractable. To circumvent this, we introduce a tractable surrogate
density $q(\bm\theta)$ from a family of distributions $\mathcal{Q}$,
and shift our focus to reformulating $\log p(\bm y)$ as an expectation
under $q(\bm\theta)$. This converts the inference problem into an
optimisation problem: instead of computing $p(\bm\theta\mid\bm y)$
directly, we search for the $q(\bm\theta)\in\mathcal{Q}$ that best
approximates it. This is made possible by noting that $\log p(\bm y)$
contains no $\bm\theta$, so it is constant with respect to~$\bm\theta$
and its value is unchanged by multiplication with any expression
equal to~1. Since every probability density integrates to~1, both
$\displaystyle\int p(\bm\theta)\,d\bm\theta$ and $\displaystyle\int q(\bm\theta)\,d\bm\theta$
are valid choices, giving}}
\log p(\bm y)
&= \log p(\bm y)\int p(\bm\theta)\,d\bm\theta \\
&= \log p(\bm y)\int q(\bm\theta)\,d\bm\theta \\
&= \int q(\bm\theta)\log p(\bm y)\,d\bm\theta \\
&= \int q(\bm\theta)f(\bm\theta)\,d\bm\theta,
\quad f(\bm\theta)\equiv \log p(\bm y) \\
&= \E_q\!\left[f(\bm\theta)\right] \\
\end{align} 
\begin{align} 
\log p(\bm y)
&= \E_q\!\left[\log p(\bm y)\right]  \\ 
&= \E_q\!\left[\log \frac{p(\bm y,\bm\theta)}{p(\bm\theta\mid\bm y)}\right] \\
&= \E_q\!\left[\log \left(\frac{p(\bm y,\bm\theta)}{q(\bm\theta)}\cdot
\frac{q(\bm\theta)}{p(\bm\theta\mid\bm y)}\right)\right] \\
&= \E_q\!\left[\log \frac{p(\bm y,\bm\theta)}{q(\bm\theta)}\right]
 +\E_q\!\left[\log \frac{q(\bm\theta)}{p(\bm\theta\mid\bm y)}\right] \\
&= \E_q[\log p(\bm y,\bm\theta)]-\E_q[\log q(\bm\theta)]
 +\E_q\!\left[\log \frac{q(\bm\theta)}{p(\bm\theta\mid\bm y)}\right].
\end{align} 

\textbf{Sequence of events (plain-language):}
\begin{enumerate}
\item Start from $\log p(\bm y)$, which is constant with respect to $\bm\theta$.
\item Multiply by $1$ using normalisation of $q$: $\int q(\bm\theta)\,d\bm\theta=1$.
\item Rewrite as an integral: $\log p(\bm y)=\int q(\bm\theta)\log p(\bm y)\,d\bm\theta$.
\item Recognise expectation form: $\int q(\bm\theta)f(\bm\theta)\,d\bm\theta=\E_q[f(\bm\theta)]$.
\item Use Bayes identity inside the logarithm: $p(\bm y)=\frac{p(\bm y,\bm\theta)}{p(\bm\theta\mid\bm y)}$.
\item Insert $q(\bm\theta)/q(\bm\theta)$ to create ELBO and KL terms.
\item Split product and quotient with log rules, then split expectations linearly.
\item Identify
\[\ELBO(q)=\E_q[\log p(\bm y,\bm\theta)]-\E_q[\log q(\bm\theta)]\]
and
\[\mathrm{KL}\!\left(q(\bm\theta)\,\|\,p(\bm\theta\mid\bm y)\right)=\E_q\!\left[\log\frac{q(\bm\theta)}{p(\bm\theta\mid\bm y)}\right].\]
\end{enumerate}

Define
\begin{align}
\ELBO(q)
&= \E_q[\log p(\bm y,\bm\theta)]-\E_q[\log q(\bm\theta)] \\
&= \E_q[\log p(\bm y,\bm\theta)] + H(q),
\end{align}
so that
\begin{equation}
\log p(\bm y) = \ELBO(q) + \mathrm{KL}\!\left(q(\bm\theta)\,\|\,p(\bm\theta\mid\bm y)\right).
\end{equation}
Since $\mathrm{KL}(\cdot\|\cdot)\ge 0$, maximising $\ELBO(q)$ is equivalent to minimising divergence to the true posterior.
 

\subsection*{Next Steps}


\section*{Alignment with Course Requirements}


\section*{Timeline}


\section*{Summary}


\end{document}
